% =============================================================
%  第 6 章 量子统计力学（Quantum statistical mechanics）
%  原文：Chapter 6, p.156--175（正文止于 p.175 上半页；p.175 下半页起为习题）
%  公式：6.1--6.93   图：6.1--6.7
% =============================================================
\bichapter{量子统计力学}{Quantum Statistical Mechanics}

经典统计力学的适用性存在局限。在低温下，纳入量子力学效应的必要性显得尤为突出。本章中，我们将首先在分子气体与固体的热容、以及黑体辐射中的紫外灾变这些情境下，展示经典结果的失效。然后我们将用量子概念重新表述统计力学。

% =============================================================
\bisection{稀薄多原子气体}{Dilute polyatomic gases}
% p.156

考虑一种稀薄的多原子分子气体。由 $n$ 个原子构成的\emph{每个}分子的 Hamilton 量是
\begin{equation}\label{eq:6.1}
\mathscr H_1=\sum_{i=1}^n\frac{\vec p_i^{\,2}}{2m}+\mathcal V(\vec q_1,\cdots,\vec q_n),
\end{equation}
其中势能 $\mathcal V$ 包含关于分子键的全部信息。为简单起见，我们假定了分子中所有原子具有相同的质量。如果质量不同，可以通过把坐标 $\vec q_i$ 按 $\sqrt{m_i/m}$ 重新标度（而把动量按 $\sqrt{m/m_i}$ 重新标度）把 Hamilton 量化成上述形式，其中 $m_i$ 是第 $i$ 个原子的质量。忽略分子\emph{之间}的相互作用，稀薄气体的配分函数为
\begin{equation}\label{eq:6.2}
Z(N)=\frac{Z_1^N}{N!}=\frac1{N!}\left\{\int\prod_{i=1}^n\frac{\mathrm d^3\vec p_i\,\mathrm d^3\vec q_i}{h^3}\exp\left[-\beta\sum_{i=1}^n\frac{\vec p_i^{\,2}}{2m}-\beta\mathcal V(\vec q_1,\cdots,\vec q_n)\right]\right\}^N.
\end{equation}
把分子结合在一起的化学键通常相当强（能量为电子伏特的量级）。在通常可达到的温度下，这些温度远小于相应的解离温度（$\approx10^4$ K），分子具有确定的形状，只发生微小的形变。这些形变对单粒子配分函数 $Z_1$ 的贡献可以如下计算：

\begin{enumerate}[label=(\alph*),leftmargin=2.2em,itemsep=0.2em]
\item 第一步是求出平衡位置 $(\vec q_1^*,\cdots,\vec q_n^*)$，即把势能 $\mathcal V$ 极小化。

% p.157

\item 围绕平衡位置的微小形变的能量代价，则通过令 $\vec q_i=\vec q_i^*+\vec u_i$、并把结果按 $\vec u$ 的幂次展开而得到，
\begin{equation}\label{eq:6.3}
\mathcal V=\mathcal V^*+\frac12\sum_{i,j=1}^n\sum_{\alpha,\beta=1}^3\frac{\partial^2\mathcal V}{\partial q_{i,\alpha}\partial q_{j,\beta}}\,u_{i,\alpha}u_{j,\beta}+\mathcal O(u^3).
\end{equation}
（这里 $i,j=1,\cdots,n$ 标记原子，而 $\alpha,\beta=1,2,3$ 标记某个特定的分量。）由于展开是围绕一个稳定的平衡构型进行的，式 \eqref{eq:6.3} 中不出现一阶导数，而二阶导数矩阵是正定的（\textit{positive definite}），也就是说，它只有非负的本征值。

\item 分子的简正模（\textit{normal modes}）通过把 $3n\times3n$ 矩阵 $\partial^2\mathcal V/\partial q_{i,\alpha}\partial q_{j,\beta}$ 对角化得到。所得的 $3n$ 个本征值 $\{K_s\}$ 表示每个模式的\term{劲度}{stiffness}{6.stiffness}。我们可以把变量从原来的形变 $\{\vec u_i\}$ 变换到本征模的振幅 $\{\tilde u_s\}$。相应的共轭动量是 $\{\tilde p_s=m\dot{\tilde u}_s\}$。由于从 $\{\vec u_i\}$ 到 $\{\tilde u_s\}$ 的变换是\term{幺正的}{unitary}{6.unitary}（保持向量的长度，$\sum_s\tilde p_s^{\,2}=\sum_i\vec p_i^{\,2}$），且所得的形变 Hamilton 量的二次型部分为
\begin{equation}\label{eq:6.4}
\mathscr H_1=\mathcal V^*+\sum_{s=1}^{3n}\left[\frac1{2m}\tilde p_s^{\,2}+\frac{K_s}{2}\tilde u_s^{\,2}\right].
\end{equation}
{\emergencystretch=2em\relax
（这类变换也是正则的（\textit{canonical}），即保持相空间中的积分测度，$\prod_{i,\alpha}\mathrm du_{i,\alpha}\mathrm dp_{i,\alpha}=\prod_s\mathrm d\tilde u_s\,\mathrm d\tilde p_s$。）\par}
\end{enumerate}

每个分子的平均能量是上述 Hamilton 量的期望值。由于每个二次自由度对能量的贡献经典地是一个 $k_BT/2$ 的因子，
\begin{equation}\label{eq:6.5}
\langle\mathscr H_1\rangle=\mathcal V^*+\frac{3n+m}{2}k_BT.
\end{equation}
只有具有有限劲度的模式才能储存势能，而 $m$ 被定义为具有非零 $K_s$ 的这种模式的数目。势的下列对称性迫使某些本征值为零：

\begin{enumerate}[label=(\alph*),leftmargin=2.2em,itemsep=0.2em]
\item \emph{平移对称性}（\textit{Translation symmetry}）：由于 $\mathcal V(\vec q_1+\vec c,\cdots,\vec q_n+\vec c)=\mathcal V(\vec q_1,\cdots,\vec q_n)$，质心坐标 $\vec Q=\sum_\alpha\vec q_\alpha/n$ 中不储存能量，即 $\mathcal V(\vec Q)=\mathcal V(\vec Q+\vec c)$，而相应的三个 $K_{\text{trans}}$ 值为零。
\item \term{旋转对称性}{Rotation symmetry}{6.rotation-symmetry}：与分子的转动相联系的势能同样不存在，相应的劲度 $K_{\text{rot}}=0$。转动模式的数目 $0\le r\le3$ 取决于分子的形状。一般 $r=3$，例外是只有一个原子的情形，它有 $r=0$（不可能有转动），以及任何棒状分子，它有 $r=2$，因为绕其轴线的转动不导致新的构型。
\end{enumerate}

矩阵的其余 $m=3n-3-r$ 个本征向量具有非零劲度，并对应于\term{振动的}{vibrational}{6.vibrational}简正模。由式 \eqref{eq:6.5}，每个分子的能量为
\begin{equation}\label{eq:6.6}
\langle\mathscr H_1\rangle=\frac{6n-3-r}{2}k_BT.
\end{equation}

% p.158

相应的热容为
\begin{equation}\label{eq:6.7}
C_V=\frac{(6n-3-r)}{2}k_B,\qquad\text{和}\qquad C_P=C_V+k_B=\frac{(6n-1-r)}{2}k_B,
\end{equation}
它们与温度无关。比值 $\gamma=C_P/C_V$ 在绝热过程中容易测量。基于上述论证所预期的 $\gamma$ 值，对若干不同的分子列在下面。

\begin{center}
\begin{tabular}{lllll}
单原子 & He & $n=1$ & $r=0$ & $\gamma=5/3$\\
双原子 & O$_2$ 或 CO & $n=2$ & $r=2$ & $\gamma=9/7$\\
线性三原子 & O--C--O & $n=3$ & $r=2$ & $\gamma=15/13$\\
平面三原子 & H$_2$O & $n=3$ & $r=3$ & $\gamma=14/12=7/6$\\
四原子 & NH$_3$ & $n=4$ & $r=3$ & $\gamma=20/18=10/9$
\end{tabular}
\end{center}

\cnfigure{fig6-1}{双原子分子气体的典型热容对温度的依赖关系。}

对稀薄气体热容的测量结果与上述预言并不一致。例如，对诸如氧气这样的双原子气体，数值 $C_V/k_B=7/2$ 只有在温度高于几千开尔文时才被观察到。在室温下观察到的是较低的数值 $5/2$，而在约 10 K 的更低的温度下，它进一步减小到 $3/2$。低温下的数值与单原子气体的数值相似，这提示转动与振动自由度中都没有储存能量。如果把允许的能级\emph{量子化}，这些观测结果就可以得到解释，如下所述。

\emph{振动模式}：双原子分子有一个振动模式，其劲度为 $K\equiv m\omega^2$，其中 $\omega$ 是振荡的频率。这个模式的经典配分函数是
\begin{equation}\label{eq:6.8}
\begin{aligned}
Z_{\text{vib}}^c&=\int\frac{\mathrm dp\,\mathrm dq}{h}\exp\left[-\beta\left(\frac{p^2}{2m}+\frac{m\omega^2q^2}{2}\right)\right]\\
&=\frac1h\sqrt{\left(\frac{2\pi m}{\beta}\right)\left(\frac{2\pi}{\beta m\omega^2}\right)}=\frac{2\pi}{h\beta\omega}=\frac{k_BT}{\hbar\omega},
\end{aligned}
\end{equation}

% p.159

其中 $\hbar=h/2\pi$。这个模式中储存的相应能量为
\begin{equation}\label{eq:6.9}
\langle\mathscr H_{\text{vib}}\rangle^c=-\frac{\partial\ln Z}{\partial\beta}=\frac{\partial\ln(\beta\hbar\omega)}{\partial\beta}=\frac1\beta=k_BT,
\end{equation}
它来自每个动能与势能自由度的 $k_BT/2$。在量子力学中，允许的能量值被量子化为
\begin{equation}\label{eq:6.10}
\mathscr H^q_{\text{vib}}=\hbar\omega\left(n+\frac12\right),
\end{equation}
其中 $n=0,1,2,\cdots$。假定每个分立能级的概率正比于它的 Boltzmann 权重（稍后将对这一点加以论证），则有一个归一化因子
\begin{equation}\label{eq:6.11}
Z^q_{\text{vib}}=\sum_{n=0}^{\infty}\mathrm e^{-\beta\hbar\omega(n+1/2)}=\frac{\mathrm e^{-\beta\hbar\omega/2}}{1-\mathrm e^{-\beta\hbar\omega}}.
\end{equation}
高温极限
\begin{equation*}
\lim_{\beta\to0}Z^q_{\text{vib}}=\frac1{\beta\hbar\omega}=\frac{k_BT}{\hbar\omega},
\end{equation*}
与式 \eqref{eq:6.8} 一致（部分原因在于选择 $h$ 作为经典相空间的测度）。

振动能量的期望值为
\begin{equation}\label{eq:6.12}
E^q_{\text{vib}}=-\frac{\partial\ln Z}{\partial\beta}=\frac{\hbar\omega}{2}+\frac{\partial}{\partial\beta}\ln\left(1-\mathrm e^{-\beta\hbar\omega}\right)=\frac{\hbar\omega}{2}+\hbar\omega\frac{\mathrm e^{-\beta\hbar\omega}}{1-\mathrm e^{-\beta\hbar\omega}}.
\end{equation}
第一项是即使在零温基态中也存在的、由\term{量子涨落}{quantum fluctuations}{6.quantum-fluctuations}引起的能量。第二项描述了由热涨落引起的附加能量。由此得到的热容为
\begin{equation}\label{eq:6.13}
C^q_{\text{vib}}=\frac{\mathrm dE^q_{\text{vib}}}{\mathrm dT}=k_B\left(\frac{\hbar\omega}{k_BT}\right)^2\frac{\mathrm e^{-\beta\hbar\omega}}{\left(1-\mathrm e^{-\beta\hbar\omega}\right)^2},
\end{equation}

只有在温度 $T\gg\theta_{\text{vib}}$ 时才达到经典值 $k_B$，其中 $\theta_{\text{vib}}=\hbar\omega/k_B$ 是与振动的量子的能量相联系的特征温度。

\cnfigure{fig6-2}{（量子化的）振动对热容的贡献。}

% p.160

对于 $T\ll\theta_{\text{vib}}$，$C^q_{\text{vib}}$ 按 $\exp(-\theta_{\text{vib}}/T)$ 趋于零。$\theta_{\text{vib}}$ 的典型值在 $10^3$ 到 $10^4$ K 的范围内，这解释了为什么只有在较高的温度下才观察到热容的经典值。

\emph{转动模式}：为了解释双原子分子热容中的低温反常，我们必须研究转动自由度的量子化。经典地，双原子分子的取向由两个角 $\theta$ 和 $\phi$ 确定，它的 Lagrange 量（等于动能）为
\begin{equation}\label{eq:6.14}
\mathscr L=\frac I2\left(\dot\theta^2+\sin^2\theta\,\dot\phi^2\right),
\end{equation}
其中 $I$ 是转动惯量。用共轭动量表示，
\begin{equation}\label{eq:6.15}
p_\theta=\frac{\partial\mathscr L}{\partial\dot\theta}=I\dot\theta,\qquad p_\phi=\frac{\partial\mathscr L}{\partial\dot\phi}=I\sin^2\theta\,\dot\phi,
\end{equation}
转动的 Hamilton 量为
\begin{equation}\label{eq:6.16}
\mathscr H_{\text{rot}}=\frac1{2I}\left(p_\theta^2+\frac{p_\phi^2}{\sin^2\theta}\right)\equiv\frac{\vec L^2}{2I},
\end{equation}
其中 $\vec L$ 是角动量。由经典配分函数，
\begin{equation}\label{eq:6.17}
\begin{aligned}
Z^c_{\text{rot}}&=\frac1{h^2}\int_0^\pi\mathrm d\theta\int_0^{2\pi}\mathrm d\phi\int_{-\infty}^{\infty}\mathrm dp_\theta\,\mathrm dp_\phi\exp\left[-\frac\beta{2I}\left(p_\theta^2+\frac{p_\phi^2}{\sin^2\theta}\right)\right]\\
&=\left(\frac{2\pi I}{\beta}\right)\left(\frac{4\pi}{h^2}\right)=\frac{2Ik_BT}{\hbar^2},
\end{aligned}
\end{equation}
储存的能量为
\begin{equation}\label{eq:6.18}
\langle E_{\text{rot}}\rangle^c=-\frac{\partial\ln Z}{\partial\beta}=\frac{\partial}{\partial\beta}\ln\left(\frac{\beta\hbar^2}{2I}\right)=k_BT,
\end{equation}
正如对两个自由度所预期的那样。在量子力学中，角动量的允许值被量子化为 $\vec L^2=\hbar^2\ell(\ell+1)$，其中 $\ell=0,1,2,\cdots$，而每个态有 $2\ell+1$ 的重数（沿某个选定的方向，$L_z=-\ell,\cdots,+\ell$）。于是对这些能级得到配分函数
\begin{equation}\label{eq:6.19}
Z^q_{\text{rot}}=\sum_{\ell=0}^{\infty}\exp\left[-\frac{\beta\hbar^2\ell(\ell+1)}{2I}\right](2\ell+1)=\sum_{\ell=0}^{\infty}\exp\left[-\frac{\theta_{\text{rot}}\,\ell(\ell+1)}{T}\right](2\ell+1),
\end{equation}
其中 $\theta_{\text{rot}}=\hbar^2/(2Ik_B)$ 是与转动能量的量子相联系的特征温度。虽然这个和一般不能解析地求出，但我们可以研究它的高温与低温极限。

\begin{enumerate}[label=(\alph*),leftmargin=2.2em,itemsep=0.2em]
\item 对于 $T\gg\theta_{\text{rot}}$，式 \eqref{eq:6.19} 中的各项变化缓慢，求和可以用积分代替
\begin{equation}\label{eq:6.20}
\begin{aligned}
\lim_{T\to\infty}Z^q_{\text{rot}}&=\int_0^{\infty}\mathrm dx(2x+1)\exp\left[-\frac{\theta_{\text{rot}}x(x+1)}{T}\right]\\
&=\int_0^{\infty}\mathrm dy\,\mathrm e^{-\theta_{\text{rot}}y/T}=\frac{T}{\theta_{\text{rot}}}=Z^c_{\text{rot}},
\end{aligned}
\end{equation}
也就是说，重新得到了式 \eqref{eq:6.17} 的经典结果。

% p.161

\item 对于 $T\ll\theta_{\text{rot}}$，求和由前几项主导，而
\begin{equation}\label{eq:6.21}
\lim_{T\to\infty}Z^q_{\text{rot}}=1+3\mathrm e^{-2\theta_{\text{rot}}/T}+\mathcal O(\mathrm e^{-6\theta_{\text{rot}}/T}),
\end{equation}
由此得到的能量为
\begin{equation}\label{eq:6.22}
E^q_{\text{rot}}=-\frac{\partial\ln Z}{\partial\beta}\approx-\frac{\partial}{\partial\beta}\ln\left[1+3\mathrm e^{-2\theta_{\text{rot}}/T}\right]\approx6k_B\theta_{\text{rot}}\,\mathrm e^{-2\theta_{\text{rot}}/T}.
\end{equation}
由此得到的热容在低温下趋于零，即
\begin{equation}\label{eq:6.23}
C_{\text{rot}}=\frac{\mathrm dE^q_{\text{rot}}}{\mathrm dT}=3k_B\left(\frac{2\theta_{\text{rot}}}{T}\right)^2\mathrm e^{-2\theta_{\text{rot}}/T}+\cdots.
\end{equation}
\end{enumerate}

\cnfigure{fig6-3}{（量子化的）转动对棒状分子热容的贡献。}

$\theta_{\text{rot}}$ 的典型值在 1 K 到 10 K 之间，这解释了热容测量中较低温度处的肩状结构。在很低的温度下，唯一的贡献来自质心的动能，分子表现得像一个单原子粒子。（由于量子统计，热容在更低的温度下还会趋于零，这一点将在讨论全同粒子时加以讨论。）

% =============================================================
\bisection{固体的振动}{Vibrations of a solid}

粒子之间的吸引相互作用首先导致低温下从气体到液体的凝聚，并最终在更低的温度下导致冻结成固体。为了讨论它的热力学，固体可以被看作一个非常大的分子，其 Hamilton 量与式 \eqref{eq:6.1} 类似，其中有 $n=N\gg1$ 个原子。于是我们可以按上一节中列出的步骤进行。

（a）固体的经典基态构型通过极小化势能 $\mathcal V$ 得到。在几乎所有的情形下，最低能量对应于原子的一种周期性排列，形成一个\term{晶格}{lattice}{6.lattice}。用三个\term{基矢}{basis vectors}{6.basis-vectors} $\hat a$、$\hat b$ 和 $\hat c$ 表示，简单晶体中原子的位置由下式给出
\begin{equation}\label{eq:6.24}
\vec q^*(\ell,m,n)=\left[\ell\hat a+m\hat b+n\hat c\right]\equiv\vec r,
\end{equation}
其中 $\{\ell,m,n\}$ 是一组整数。

% p.162

（b）在有限的温度下，原子可以发生微小的形变
\begin{equation}\label{eq:6.25}
\vec q_{\vec r}=\vec r+\vec u(\vec r),
\end{equation}
其势能代价为
\begin{equation}\label{eq:6.26}
\mathcal V=\mathcal V^*+\frac12\sum_{\vec r,\vec r'}\sum_{\alpha,\beta}\frac{\partial^2\mathcal V}{\partial q_{\vec r,\alpha}\partial q_{\vec r',\beta}}u_\alpha(\vec r)u_\beta(\vec r')+\mathcal O(u^3).
\end{equation}

（c）找出晶体的简正模，因其\emph{平移对称性}（\textit{translational symmetry}）而大大简化。特别地，二阶导数矩阵只依赖于两点之间的\emph{相对的}（\textit{relative}）距离，
\begin{equation}\label{eq:6.27}
\frac{\partial^2\mathcal V}{\partial q_{\vec r,\alpha}\partial q_{\vec r',\beta}}=K_{\alpha\beta}(\vec r-\vec r').
\end{equation}
总是可以利用这样的对称性，通过采用一个\term{Fourier 基}{Fourier basis}{6.fourier-basis}来至少部分地把这个矩阵对角化，
\begin{equation}\label{eq:6.28}
u_\alpha(\vec r)=\sum_{\vec k}{}'\frac{\mathrm e^{\mathrm i\vec k\cdot\vec r}}{\sqrt N}\tilde u_\alpha(\vec k).
\end{equation}
这个求和限定在\term{布里渊区}{Brillouin zone}{6.brillouin-zone}内部的\term{波矢}{wavevectors}{6.wavevector} $\vec k$ 上。例如，在间距为 $a$ 的立方晶格中，$\vec k$ 的每个分量被限制在区间 $[-\pi/a,\pi/a]$ 内。这是因为该区间之外的波矢不携带任何附加信息，因为 $(k_x+2\pi m/a)(na)=k_x(na)+2mn\pi$，而相差 $2\pi$ 整数倍的相位不影响式 \eqref{eq:6.28} 中的求和。

\cnfigure{fig6-4}{间距为 $a$ 的正方晶格的布里渊区。}

用 Fourier 模表示，形变的势能为
\begin{equation}\label{eq:6.29}
\mathcal V=\mathcal V^*+\frac1{2N}\sum_{(\vec r,\vec r'),(\vec k,\vec k')}\sum_{\alpha,\beta}K_{\alpha\beta}(\vec r-\vec r')\,\mathrm e^{\mathrm i\vec k\cdot\vec r}\tilde u_\alpha(\vec k)\,\mathrm e^{\mathrm i\vec k'\cdot\vec r'}\tilde u_\beta(\vec k').
\end{equation}
我们可以把变量变换到相对坐标与质心坐标，
\begin{equation*}
\vec\rho=\vec r-\vec r',\qquad\text{和}\qquad\vec R=\frac{\vec r+\vec r'}2,
\end{equation*}

% p.163

即令
\begin{equation*}
\vec r=\vec R+\frac{\vec\rho}2,\qquad\text{和}\qquad\vec r'=\vec R-\frac{\vec\rho}2.
\end{equation*}
式 \eqref{eq:6.29} 现在简化为
\begin{equation}\label{eq:6.30}
\mathcal V=\mathcal V^*+\frac1{2N}\sum_{\vec k,\vec k'}\left(\sum_{\vec R}\mathrm e^{\mathrm i(\vec k+\vec k')\cdot\vec R}\right)\left(\sum_{\vec\rho,\alpha,\beta}K_{\alpha\beta}(\vec\rho)\,\mathrm e^{\mathrm i(\vec k-\vec k')\cdot\vec\rho/2}\tilde u_\alpha(\vec k)\tilde u_\beta(\vec k')\right).
\end{equation}
由于第一个括号中的求和是 $N\delta_{\vec k+\vec k',0}$，
\begin{equation}\label{eq:6.31}
\begin{aligned}
\mathcal V&=\mathcal V^*+\frac12\sum_{\vec k,\alpha,\beta}\left[\sum_{\vec\rho}K_{\alpha\beta}(\vec\rho)\,\mathrm e^{\mathrm i\vec k\cdot\vec\rho}\right]\tilde u_\alpha(\vec k)\tilde u_\beta(-\vec k)\\
&=\mathcal V^*+\frac12\sum_{\vec k,\alpha,\beta}\tilde K_{\alpha\beta}(\vec k)\tilde u_\alpha(\vec k)\tilde u_\beta(\vec k)^*,
\end{aligned}
\end{equation}
其中 $\tilde K_{\alpha\beta}(\vec k)=\sum_{\vec\rho}K_{\alpha\beta}(\vec\rho)\exp(\mathrm i\vec k\cdot\vec\rho)$，而 $\tilde u_\beta(\vec k)^*=\tilde u_\beta(-\vec k)$ 是 $\tilde u_\beta(\vec k)$ 的复共轭。

于是，不同的 Fourier 模在二次阶上解耦，而对角化 $3N\times3N$ 的二阶导数矩阵的任务被简化为对每个 $\vec k$ 分别对角化 $3\times3$ 矩阵 $\tilde K_{\alpha\beta}(\vec k)$。$\tilde K_{\alpha\beta}$ 的形式进一步受晶体点群对称性的限制。对这种约束的讨论超出了本节的意图，为简单起见，我们将假定 $\tilde K_{\alpha\beta}(\vec k)=\delta_{\alpha,\beta}\tilde K(\vec k)$ 已经是对角化的。（对一种各向同性的材料，这意味着体模量与剪切模量之间存在一个特定的关系。）

形变的动能为
\begin{equation}\label{eq:6.32}
\sum_{i=1}^N\frac m2\dot{\vec q}_i^{\,2}=\sum_{\vec k,\alpha}\frac m2\dot{\tilde u}_\alpha(\vec k)\dot{\tilde u}_\alpha(\vec k)^*=\sum_{\vec k,\alpha}\frac1{2m}\tilde p_\alpha(\vec k)\tilde p_\alpha(\vec k)^*,
\end{equation}
其中
\begin{equation*}
\tilde p_\alpha(\vec k)=\frac{\partial\mathscr L}{\partial\dot{\tilde u}_\alpha(\vec k)}=m\dot{\tilde u}_\alpha(\vec k)
\end{equation*}
是 $\tilde u_\alpha(\vec k)$ 的共轭动量。所得的形变 Hamilton 量
\begin{equation}\label{eq:6.33}
\mathscr H=\mathcal V^*+\sum_{\vec k,\alpha}\left[\frac1{2m}\left|\tilde p_\alpha(\vec k)\right|^2+\frac{\tilde K(\vec k)}{2}\left|\tilde u_\alpha(\vec k)\right|^2\right],
\end{equation}
描述了 $3N$ 个独立的谐振子，其频率为 $\omega_\alpha(\vec k)=\sqrt{\tilde K(\vec k)/m}$。

在经典处理中，每个具有非零劲度的谐振子对内能的贡献是 $k_BT$。

% p.164

晶体的转动）。因此，在相差 $1/N$ 量级的非广延修正的意义上，与 Hamilton 量 (6.33) 相联系的内能经典地是 $3Nk_BT$，从而给出每个原子的热容为 $3k_B$，且与温度无关。事实上，测量到的热容在低温下趋于零。我们同样可以把这一观测结果与每个振子的能级的量子化联系起来，如上一节所讨论的那样。分别把每个简谐模量子化，得到 Hamilton 量的允许值为
\begin{equation}\label{eq:6.34}
\mathscr H^q=\mathcal V^*+\sum_{\vec k,\alpha}\hbar\omega_\alpha(\vec k)\left(n_{\vec k,\alpha}+\frac12\right),
\end{equation}
其中整数集合 $\{n_{\vec k,\alpha}\}$ 描述了振子的量子微观态。由于这些振子是独立的，它们的配分函数
\begin{equation}\label{eq:6.35}
Z^q=\sum_{\{n_{\vec k,\alpha}\}}\mathrm e^{-\beta\mathscr H^q}=\mathrm e^{-\beta E_0}\prod_{\vec k,\alpha}\sum_{n_{\vec k,\alpha}}\mathrm e^{-\beta\hbar\omega_{\vec k,\alpha}n_{\vec k,\alpha}}=\mathrm e^{-\beta E_0}\prod_{\vec k,\alpha}\left[\frac1{1-\mathrm e^{-\beta\hbar\omega_{\vec k,\alpha}}}\right]
\end{equation}
是单振子配分函数（如式 \eqref{eq:6.11}）的乘积。（$E_0$ 除了包含 $\mathcal V^*$ 之外，还包含所有振子的基态能量。）

内能为
\begin{equation}\label{eq:6.36}
E(T)=\langle\mathscr H^q\rangle=E_0+\sum_{\vec k,\alpha}\hbar\omega_\alpha(\vec k)\left\langle n_\alpha(\vec k)\right\rangle,
\end{equation}
其中平均\emph{占有数}（\textit{occupation numbers}）由下式给出
\begin{equation}\label{eq:6.37}
\begin{aligned}
\left\langle n_\alpha(\vec k)\right\rangle&=\frac{\sum_{n=0}^\infty n\,\mathrm e^{-\beta\hbar\omega_\alpha(\vec k)n}}{\sum_{n=0}^\infty\mathrm e^{-\beta\hbar\omega_\alpha(\vec k)n}}=-\frac{\partial}{\partial(\beta\hbar\omega_\alpha(\vec k))}\ln\left(\frac1{1-\mathrm e^{-\beta\hbar\omega_\alpha(\vec k)}}\right)\\
&=\frac{\mathrm e^{-\beta\hbar\omega_\alpha(\vec k)}}{1-\mathrm e^{-\beta\hbar\omega_\alpha(\vec k)}}=\frac1{\mathrm e^{\beta\hbar\omega_\alpha(\vec k)}-1}.
\end{aligned}
\end{equation}

作为纳入量子力学效应的第一次尝试，我们可以采用\term{爱因斯坦模型}{Einstein model}{6.einstein-model}，其中假定所有振子具有相同的频率 $\omega_E$。这个模型对应于原子被劲度为 $\tilde K=\partial^2\mathcal V/\partial q^2=m\omega_E^2$ 的弹簧固定在其理想位置上。由此得到的内能
\begin{equation}\label{eq:6.38}
E=E_0+3N\frac{\hbar\omega_E\,\mathrm e^{-\beta\hbar\omega_E}}{1-\mathrm e^{-\beta\hbar\omega_E}},
\end{equation}
与热容
\begin{equation}\label{eq:6.39}
C=\frac{\mathrm dE}{\mathrm dT}=3Nk_B\left(\frac{T_E}{T}\right)^2\frac{\mathrm e^{-T_E/T}}{\left(1-\mathrm e^{-T_E/T}\right)^2},
\end{equation}
都简单地正比于单个振子的相应量（式 \eqref{eq:6.12} 与 \eqref{eq:6.13}）。特别地，热容以特征温度 $T_E=\hbar\omega_E/k_B$ 指数式地衰减到零。然而，实验测量到的热容衰减到零要慢得多，为 $T^3$。

这个分歧通过\term{德拜模型}{Debye model}{6.debye-model}得到解决，它强调在低温下对热容的主要贡献来自

% p.165

\cnfigure{fig6-5}{固体的热容（灰线）比在爱因斯坦模型中（实线）更缓慢地趋于零。}

频率最低、也就是最容易被激发的那些振子。能量最低的模式又对应于最小的波矢 $k=|\vec k|$，或者说最长的波长 $\lambda=2\pi/k$。确实，$\vec k=0$ 的模式只不过描述了晶格的纯平移，其劲度为零。由连续性，我们预期 $\lim_{\vec k\to0}\tilde K(\vec k)=0$，而忽略晶体对称性的考虑，$\tilde K(\vec k)$ 在小波矢处的展开具有形式
\begin{equation*}
\tilde K(\vec k)=Bk^2+\mathcal O(k^4).
\end{equation*}
展开式中不出现奇次项，因为在实空间中 $K(\vec r-\vec r')=K(\vec r'-\vec r)$ 蕴含 $\tilde K(\vec k)=\tilde K(-\vec k)$。长波长处的相应频率为
\begin{equation}\label{eq:6.40}
\omega(\vec k)=\sqrt{\frac{Bk^2}{m}}=vk,
\end{equation}
其中 $v=\sqrt{B/m}$ 是晶体中的\term{声速}{speed of sound}{6.speed-of-sound}。（在真实的材料中，$\tilde K_{\alpha\beta}$ 不正比于 $\delta_{\alpha\beta}$，不同\emph{偏振}（\textit{polarizations}）的声波具有不同的速度。）

\cnfigure{fig6-6}{各向同性晶体中声子的色散对纵波与横波有不同的分支。}

% p.166

振动模式的量子通常称为\term{声子}{phonons}{6.phonons}。利用式 \eqref{eq:6.40} 中的\term{色散关系}{dispersion relation}{6.dispersion-relation}，低能声子对内能的贡献为
\begin{equation}\label{eq:6.41}
\langle\mathscr H^q\rangle=E_0+\sum_{\vec k,\alpha}\frac{\hbar vk}{\mathrm e^{\beta\hbar vk}-1}.
\end{equation}
在尺寸为 $L_x\times L_y\times L_z$ 的盒子中取\term{周期性边界条件}{periodic boundary conditions}{6.periodic-bc}，允许的波矢为
\begin{equation}\label{eq:6.42}
\vec k=\left(\frac{2\pi n_x}{L_x},\frac{2\pi n_y}{L_y},\frac{2\pi n_z}{L_z}\right),
\end{equation}
其中 $n_x$、$n_y$、$n_z$ 是整数。在尺寸很大的极限下，这些模式分布得非常密集，而体积元 $\mathrm d^3\vec k$ 中的模式数为
\begin{equation}\label{eq:6.43}
\mathrm dN=\frac{\mathrm dk_x}{2\pi/L_x}\frac{\mathrm dk_y}{2\pi/L_y}\frac{\mathrm dk_z}{2\pi/L_z}=\frac V{(2\pi)^3}\mathrm d^3\vec k\equiv\rho\,\mathrm d^3\vec k.
\end{equation}
利用\term{态密度}{density of states}{6.density-of-states} $\rho$，任何对允许波矢的求和都可以用一个积分代替，即
\begin{equation}\label{eq:6.44}
\lim_{V\to\infty}\sum_{\vec k}f(\vec k)=\int\mathrm d^3\vec k\,\rho f(\vec k).
\end{equation}
因此，式 \eqref{eq:6.41} 可以改写为
\begin{equation}\label{eq:6.45}
E=E_0+3V\int^{\text{B.Z.}}\frac{\mathrm d^3\vec k}{(2\pi)^3}\frac{\hbar vk}{\mathrm e^{\beta\hbar vk}-1},
\end{equation}
其中积分在布里渊区的体积上进行，而因子 3 来自假定三个偏振具有相同的声速。

由于它依赖于布里渊区的形状，我们无法对式 \eqref{eq:6.45} 中的能量给出一个简单的闭式表达式。不过，我们可以考察它的高温与低温极限。把这两个极限分开的特征\term{德拜温度}{Debye temperature}{6.debye-temperature}
\begin{equation}\label{eq:6.46}
T_D=\frac{\hbar vk_{\text{max}}}{k_B}\approx\frac{\hbar v}{k_B}\cdot\frac\pi a,
\end{equation}
对应于布里渊区边缘处的高频模式。对于 $T\gg T_D$，所有模式都表现为经典行为。式 \eqref{eq:6.45} 的被积函数就是 $k_BT$，而由于模式的总数是 $3N=3V\int^{\text{B.Z.}}\mathrm d^3\vec k/(2\pi)^3$，就重新得到 $E(T)=E_0+3Nk_BT$ 与 $C=3Nk_B$ 这两个经典结果。对于 $T\ll T_D$，式 \eqref{eq:6.45} 分母中的因子 $\exp(\beta\hbar vk)$ 在布里渊区边缘处非常大。积分最重要的贡献来自小的 $k$，而把积分范围扩展到无穷所引起的误差很小。在把

% p.167

变量变换为 $x=\beta\hbar v|\vec k|$、并由于球对称性而利用 $\mathrm d^3\vec k=4\pi x^2\,\mathrm dx/(\beta\hbar v)^3$ 之后，式 \eqref{eq:6.45} 给出
\begin{equation}\label{eq:6.47}
\begin{aligned}
\lim_{T\ll T_D}E(T)&\approx E_0+\frac{3V}{8\pi^3}\left(\frac{k_BT}{\hbar v}\right)^3 4\pi k_BT\int_0^\infty\mathrm dx\frac{x^3}{\mathrm e^x-1}\\
&=E_0+\frac{\pi^2}{10}V\left(\frac{kT}{\hbar v}\right)^3k_BT.
\end{aligned}
\end{equation}
（定积分 $\pi^4/15\approx6.5$ 的数值可以在标准表中查到。）由此得到的热容
\begin{equation}\label{eq:6.48}
C=\frac{\mathrm dE}{\mathrm dT}=k_BV\frac{2\pi^2}{5}\left(\frac{k_BT}{\hbar v}\right)^3
\end{equation}
具有 $C\propto Nk_B(T/T_D)^3$ 的形式，与观测结果一致。这一结果的物理解释如下：在温度 $T\ll T_D$ 下，只有一部分声子模式能够被热激发。这些是能量量子 $\hbar\omega(\vec k)\ll k_BT$ 的低频声子。被激发的声子具有波矢 $|\vec k|<k^*(T)\approx(k_BT/\hbar v)$。相当一般地，在 $d$ 维空间中，这些模式的数目近似为 $Vk^*(T)^d\sim V(k_BT/\hbar v)^d$。每个被激发的模式都可以按经典处理，并对内能贡献大约 $k_BT$，于是内能标度为 $E\sim V(k_BT/\hbar v)^dk_BT$。相应的热容 $C\sim Vk_B(k_BT/\hbar v)^d$ 按 $T^d$ 趋于零。

% =============================================================
\bisection{黑体辐射}{Black-body radiation}

声子对应于固体介质的振动。在液态与气态中也存在（纵）声模式。然而，即使是「空」真空也能支持电磁（EM）场的涨落，即\term{光子}{photons}{6.photons}，它们在有限温度下可以被热激发。这个场的简正模是 EM 波，由波数 $\vec k$ 和两个可能的\emph{偏振}（\textit{polarizations}）$\alpha$ 表征。（由于在自由空间中 $\nabla\cdot\vec E=0$，电场必须垂直于 $\vec k$，因此只存在\term{横模}{transverse}{6.transverse}。）适当选择坐标后，EM 场的 Hamilton 量可以写成谐振子之和，
\begin{equation}\label{eq:6.49}
\mathscr H=\frac12\sum_{\vec k,\alpha}\left[\left|\tilde p_{\vec k,\alpha}\right|^2+\omega_\alpha(\vec k)^2\left|\tilde u_\alpha(\vec k)\right|^2\right],
\end{equation}
其中 $\omega_\alpha(\vec k)=ck$，$c$ 是光速。

{\emergencystretch=2em\relax
在周期性边界条件下，尺寸为 $L$ 的盒子中允许的波矢是 $\vec k=2\pi(n_x,n_y,n_z)/L$，其中 $\{n_x,n_y,n_z\}$ 是整数。然而，与声子不同，不存在限制 $\vec k$ 大小的布里渊区，这些整数可以任意大。缺乏这种限制在经典处理中导致\term{紫外灾变}{ultraviolet catastrophe}{6.ultraviolet-catastrophe}：由于对波矢没有限制，为每个模式指定 $k_BT$ 会导致在高频模式中储存无限大的能量。（低频被盒子的有限尺寸截断。）正是\par}

% p.168

为了解决这一困难，Planck 提出 EM 能量的允许值必须按照以下 Hamilton 量量子化
\begin{equation}\label{eq:6.50}
\mathscr H^q=\sum_{\vec k,\alpha}\hbar ck\left(n_\alpha(\vec k)+\frac12\right),\qquad\text{其中}\qquad n_\alpha(\vec k)=0,1,2,\cdots.
\end{equation}
与对声子的情形一样，内能由下式计算
\begin{equation}\label{eq:6.51}
E=\langle\mathscr H^q\rangle=\sum_{\vec k,\alpha}\hbar ck\left(\frac12+\frac{\mathrm e^{-\beta\hbar ck}}{1-\mathrm e^{-\beta\hbar ck}}\right)=VE_0+\frac{2V}{(2\pi)^3}\int\mathrm d^3\vec k\frac{\hbar ck}{\mathrm e^{\beta\hbar ck}-1}.
\end{equation}
零点能实际上是无穷大，但由于测量到的只是能量差，它通常被忽略。把变量变换为 $x=\beta\hbar ck$，使我们能够计算出激发能，
\begin{equation}\label{eq:6.52}
\begin{aligned}
\frac{E^*}{V}&=\frac{\hbar c}{\pi^2}\left(\frac{k_BT}{\hbar c}\right)^4\int_0^\infty\frac{\mathrm dx\,x^3}{\mathrm e^x-1}\\
&=\frac{\pi^2}{15}\left(\frac{k_BT}{\hbar c}\right)^3k_BT.
\end{aligned}
\end{equation}

EM 辐射也对容器的壁施加压力。由 Hamilton 量 (6.50) 的配分函数，
\begin{equation}\label{eq:6.53}
Z=\sum_{\{n_\alpha(\vec k)\}}\prod_{\vec k,\alpha}\exp\left[-\beta\hbar\omega(\vec k)\left(n_\alpha(\vec k)+\frac12\right)\right]=\prod_{\vec k,\alpha}\frac{\mathrm e^{-\beta\hbar ck/2}}{1-\mathrm e^{-\beta\hbar ck}},
\end{equation}
自由能为
\begin{equation}\label{eq:6.54}
\begin{aligned}
F&=-k_BT\ln Z=k_BT\sum_{\vec k,\alpha}\left[\frac{\beta\hbar ck}{2}+\ln\left(1-\mathrm e^{-\beta\hbar ck}\right)\right]\\
&=2V\int\frac{\mathrm d^3\vec k}{(2\pi)^3}\left[\frac{\hbar ck}{2}+k_BT\ln\left(1-\mathrm e^{-\beta\hbar ck}\right)\right].
\end{aligned}
\end{equation}
光子气体引起的压力为
\begin{equation}\label{eq:6.55}
\begin{aligned}
P&=-\left.\frac{\partial F}{\partial V}\right|_T=-\int\frac{\mathrm d^3\vec k}{(2\pi)^3}\left[\hbar ck+2k_BT\ln\left(1-\mathrm e^{-\beta\hbar ck}\right)\right]\\
&=P_0-\frac{k_BT}{\pi^2}\int_0^\infty\mathrm dk\,k^2\ln\left(1-\mathrm e^{-\beta\hbar ck}\right)\qquad\text{（分部积分）}\\
&=P_0+\frac{k_BT}{\pi^2}\int_0^\infty\mathrm dk\,\frac{k^3}{3}\frac{\beta\hbar c\,\mathrm e^{-\beta\hbar ck}}{1-\mathrm e^{-\beta\hbar ck}}\qquad\text{（与式 \eqref{eq:6.51} 比较）}\\
&=P_0+\frac13\frac EV.
\end{aligned}
\end{equation}
{\emergencystretch=2em\relax
注意这里也存在一个无穷大的零点压力 $P_0$。这一压力中的差值导致导体板之间可测量的\term{Casimir 力}{Casimir force}{6.casimir-force}。\par}

1/3 倍能量密度的额外压力可以与相对论性粒子气体的相应结果相比较。（如习题中所展示的，色散关系 $\mathcal E\propto|\vec p|^s$ 在 $d$ 维中导致压力 $P=(s/d)(E/V)$。）

% p.169

继续与粒子气体的类比，如果在容器的壁上开一个小孔，则单位面积、单位时间逃逸出去的能量通量为
\begin{equation}\label{eq:6.56}
\phi=\langle c_\perp\rangle\frac EV.
\end{equation}
所有光子都具有速率 $c$，而垂直于小孔的速度分量平均值的计算如下
\begin{equation}\label{eq:6.57}
\langle c_\perp\rangle=c\times\frac1{4\pi}\int_0^{\pi/2}2\pi\sin\theta\,\mathrm d\theta\cos\theta=\frac c4,
\end{equation}
由此得到
\begin{equation}\label{eq:6.58}
\phi=\frac14c\frac EV=\frac{\pi^2}{60}\frac{k_B^4T^4}{\hbar^3c^2}.
\end{equation}
这个结果 $\phi=\sigma T^4$ 就是黑体辐射的\term{Stefan--Boltzmann 定律}{Stefan--Boltzmann law}{6.stefan-boltzmann-law}，而
\begin{equation}\label{eq:6.59}
\sigma=\frac{\pi^2}{60}\frac{k_B^4}{\hbar^3c^2}\approx5.67\times10^{-8}\,\mathrm{W\,m^{-2}\,K^{-4}}
\end{equation}
是\term{Stefan 常量}{Stefan's constant}{6.stefan-constant}。黑体辐射有一个特征性的频率依赖关系：令 $E(T)/V=\int\mathrm dk\,\mathcal E(k,T)$，其中
\begin{equation}\label{eq:6.60}
\mathcal E(k,T)=\frac{\hbar c}{\pi^2}\frac{k^3}{\mathrm e^{\beta\hbar ck}-1}
\end{equation}
是波矢 $k$ 处的能量密度。区间 $[k,k+\mathrm dk]$ 中发射出的辐射通量为 $I(k,T)\mathrm dk$，其中
\begin{equation}\label{eq:6.61}
I(k,T)=\frac c4\mathcal E(k,T)=\frac{\hbar c^2}{4\pi^2}\frac{k^3}{\mathrm e^{\beta\hbar ck}-1}\to\begin{cases}ck_BTk^2/4\pi&\text{对于 }k\ll k^*(T)\\[2pt]\hbar c^2k^3\mathrm e^{-\beta\hbar ck}/4\pi^2&\text{对于 }k\gg k^*(T).\end{cases}
\end{equation}
特征波矢 $k^*(T)\approx k_BT/\hbar c$ 把量子与经典区域分开。它提供了消除紫外灾变的上截断。

\cnfigure{fig6-7}{温度为 $T$ 的系统发出的黑体辐射谱。}

% =============================================================
\bisection{量子微观态}{Quantum microstates}

% p.170

在前几节中，我们指出了经典统计力学的若干失效之处，并通过假定能级量子化而启发式地弥补了它们，同时仍然从配分和 $Z=\sum_n\exp(-\beta E_n)$ 计算热力学量。这隐含地假定了量子系统的微观态由它的分立能级来指定，并服从一个类似于 Boltzmann 权重的概率分布。这种与经典统计力学的「类比」需要加以论证。此外，量子力学本身内在地是概率性的，其不确定性不同于那些在统计力学中导致概率的不确定性。这里，我们将通过密切遵循导致经典表述的那些步骤，来构造统计力学的量子表述。

粒子系统的经典\emph{微观态}（\textit{Microstates}）由坐标与共轭动量的集合 $\{\vec p_i,\vec q_i\}$ 描述；也就是说，由 $6N$ 维相空间中的一个点描述。在量子力学中，$\{\vec q_i\}$ 与 $\{\vec p_i\}$ 不是独立的可观测量。取而代之的是：

\begin{itemize}[leftmargin=1.6em,itemsep=0.4em]
\item 量子系统的（微观）态由\term{单位向量}{unit vector}{6.unit-vector} $|\Psi\rangle$ 完全指定，它属于一个无穷维的\term{Hilbert 空间}{Hilbert space}{6.hilbert-space}。向量 $|\Psi\rangle$ 可以沿一组适当的正交归一\emph{基矢}（\textit{basis vectors}）$|n\rangle$ 用它的分量 $\langle n|\Psi\rangle$（这些是复数）表示出来。用 Dirac 引入的方便记号，这一分解可以写成
\begin{equation}\label{eq:6.62}
|\Psi\rangle=\sum_n\langle n|\Psi\rangle\,|n\rangle.
\end{equation}
最熟悉的基是坐标基 $\{|\vec q_i\rangle\}$，而 $\langle\vec q_i|\Psi\rangle\equiv\Psi(\vec q_1,\cdots,\vec q_N)$ 就是\term{波函数}{wave function}{6.wave-function}。归一化条件为
\begin{equation}\label{eq:6.63}
\langle\Psi|\Psi\rangle=\sum_n\langle\Psi|n\rangle\langle n|\Psi\rangle=1,\qquad\text{其中}\qquad\langle\Psi|n\rangle\equiv\langle n|\Psi\rangle^*.
\end{equation}
例如，在坐标基中我们必须要求
\begin{equation}\label{eq:6.64}
\langle\Psi|\Psi\rangle=\int\prod_{i=1}^N\mathrm d^3\vec q_i\,|\Psi(\vec q_1,\cdots,\vec q_N)|^2=1.
\end{equation}

\item 经典地，各种可观测量都是定义在相空间中的函数 $\mathcal O(\{\vec p_i,\vec q_i\})$。在量子力学中，这些函数被替换为 Hilbert 空间中的\term{厄米矩阵}{Hermitian matrices}{6.hermitian-matrix}（算符），其做法是在经典表达式中把 $\{\vec q_i\}$ 与 $\{\vec p_i\}$ 换成算符（在把乘积适当地对称化之后，例如 $pq\to(pq+qp)/2$）。这些基本算符满足\term{对易关系}{commutation relations}{6.commutation-relations}
\begin{equation}\label{eq:6.65}
\left[p_j,q_k\right]\equiv p_jq_k-q_kp_j=\frac\hbar{\mathrm i}\delta_{j,k}.
\end{equation}
例如，在坐标基 $\{|\vec q_i\rangle\}$ 中，动量算符是
\begin{equation}\label{eq:6.66}
p_j=\frac\hbar{\mathrm i}\frac{\partial}{\partial q_j}.
\end{equation}

% p.171

（注意经典的 Poisson 括号满足 $\{p_j,q_k\}=-\delta_{j,k}$。相当一般地，量子对易关系通过把相应的经典 Poisson 括号乘以 $\mathrm i\hbar$ 得到。）

\item 与经典力学中不同，算符 $\mathcal O$ 的取值对于一个特定的微观态并不是唯一确定的。它反而是一个随机变量，其在态 $|\Psi\rangle$ 中的\emph{平均}值由下式给出
\begin{equation}\label{eq:6.67}
\langle\mathcal O\rangle\equiv\langle\Psi|\mathcal O|\Psi\rangle\equiv\sum_{m,n}\langle\Psi|m\rangle\langle m|\mathcal O|n\rangle\langle n|\Psi\rangle.
\end{equation}
例如，
\begin{equation*}
\begin{aligned}
\langle U(\{\vec q\})\rangle&=\int\prod_{i=1}^N\mathrm d^3\vec q_i\,\Psi^*(\vec q_1,\cdots,\vec q_N)U(\{\vec q\})\Psi(\vec q_1,\cdots,\vec q_N),\qquad\text{以及}\\
\langle K(\{\vec p\})\rangle&=\int\prod_{i=1}^N\mathrm d^3\vec q_i\,\Psi^*(\vec q_1,\cdots,\vec q_N)K\left(\left\{\frac\hbar{\mathrm i}\frac{\mathrm d}{\mathrm d\vec q}\right\}\right)\Psi(\vec q_1,\cdots,\vec q_N).
\end{aligned}
\end{equation*}

为了保证期望值 $\langle\mathcal O\rangle$ 在所有态中都是实数，算符 $\mathcal O$ 必须是\emph{厄米的}（\textit{Hermitian}），也就是说，满足
\begin{equation}\label{eq:6.68}
\mathcal O^\dagger=\mathcal O,\qquad\text{其中}\qquad\langle m|\mathcal O^\dagger|n\rangle\equiv\langle n|\mathcal O|m\rangle^*.
\end{equation}
在把经典算符 $\mathcal O(\{\vec p_i,\vec q_i\})$ 中的 $\vec p$ 与 $\vec q$ 替换为相应的矩阵时，必须把各种乘积适当地对称化，以保证上述厄米性。

微观态的\emph{时间演化}（\textit{Time evolution}）由 Hamilton 量 $\mathscr H(\{\vec p_i,\vec q_i\})$ 支配。经典微观态按 Hamilton 运动方程演化，而在量子力学中态矢量随时间的变化为
\begin{equation}\label{eq:6.69}
\mathrm i\hbar\frac{\partial}{\partial t}|\Psi(t)\rangle=\mathscr H|\Psi(t)\rangle.
\end{equation}
一种方便的基是把矩阵 $\mathscr H$ \emph{对角化}的基。能量\term{本征态}{eigenstates}{6.eigenstates}满足 $\mathscr H|n\rangle=\mathcal E_n|n\rangle$，其中 $\mathcal E_n$ 是\term{本征能量}{eigenenergies}{6.eigenenergies}。把 $|\Psi(t)\rangle=\sum_n|n\rangle\langle n|\Psi(t)\rangle$ 代入式 \eqref{eq:6.69}，并利用正交归一条件 $\langle m|n\rangle=\delta_{m,n}$，得到
\begin{equation}\label{eq:6.70}
\mathrm i\hbar\frac{\mathrm d}{\mathrm dt}\langle n|\Psi(t)\rangle=\mathcal E_n\langle n|\Psi(t)\rangle,\qquad\Longrightarrow\qquad\langle n|\Psi(t)\rangle=\exp\left(-\frac{\mathrm i\mathcal E_nt}\hbar\right)\langle n|\Psi(0)\rangle.
\end{equation}
两个不同时刻的量子态可以通过一个\term{时间演化算符}{time evolution operator}{6.time-evolution-operator}相联系，即
\begin{equation}\label{eq:6.71}
|\Psi(t)\rangle=U(t,t_0)|\Psi(t_0)\rangle,
\end{equation}
它满足 $\mathrm i\hbar\partial_tU(t,t_0)=\mathscr H U(t,t_0)$，边界条件为 $U(t_0,t_0)=1$。如果 $\mathscr H$ 与 $t$ \emph{无关}，我们可以解出这些方程，得到
\begin{equation}\label{eq:6.72}
U(t,t_0)=\exp\left[-\frac{\mathrm i}\hbar\mathscr H(t-t_0)\right].
\end{equation}
\end{itemize}

% =============================================================
\bisection{量子宏观态}{Quantum macrostates}

% p.172

系统的\emph{宏观态}（\textit{Macrostates}）只依赖于少数几个热力学函数。对对应于一个给定宏观态，我们可以由大量 $\mathcal N$ 个微观态 $\mu_\alpha$ 组成一个\emph{系综}（\textit{ensemble}）。不同的微观态以概率 $p_\alpha$ 出现。（例如，在没有任何其它信息时，$p_\alpha=1/\mathcal N$ 给出一个\emph{无偏}（\textit{unbiased}）估计。）当我们不再精确地知道一个系统的微观态时，就说该系统处于一个\emph{混合态}（\textit{mixed state}）。

经典地，系综平均由下式计算
\begin{equation}\label{eq:6.73}
\overline{\mathcal O(\{\vec p_i,\vec q_i\})}_t=\sum_\alpha p_\alpha\,\mathcal O(\mu_\alpha(t))=\int\prod_{i=1}^N\mathrm d^3\tilde p_i\,\mathrm d^3\tilde q_i\,\mathcal O(\{\tilde p_i,\tilde q_i\})\,\rho(\{\tilde p_i,\tilde q_i\},t),
\end{equation}
其中
\begin{equation}\label{eq:6.74}
\rho(\{\tilde p_i,\tilde q_i\},t)=\sum_\alpha p_\alpha\prod_{i=1}^N\delta^3\big(\tilde q_i-\tilde q_i(t)_\alpha\big)\delta^3\big(\tilde p_i-\tilde p_i(t)_\alpha\big)
\end{equation}
是\term{相空间密度}{phase space density}{6.phase-space-density}。

类似地，混合量子态由一组可能的状态 $\{|\Psi_\alpha\rangle\}$ 得到，其概率为 $\{p_\alpha\}$。于是式 \eqref{eq:6.67} 中量子力学期望值的系综平均为
\begin{equation}\label{eq:6.75}
\begin{aligned}
\overline{\langle\mathcal O\rangle}&=\sum_\alpha p_\alpha\langle\Psi_\alpha|\mathcal O|\Psi_\alpha\rangle=\sum_{\alpha;m,n}p_\alpha\langle\Psi_\alpha|m\rangle\langle n|\Psi_\alpha\rangle\langle m|\mathcal O|n\rangle\\
&=\sum_{m,n}\langle n|\rho|m\rangle\langle m|\mathcal O|n\rangle=\operatorname{tr}(\rho\mathcal O),
\end{aligned}
\end{equation}
其中我们引入了一组基 $\{|n\rangle\}$，并定义了\term{密度矩阵}{density matrix}{6.density-matrix}
\begin{equation}\label{eq:6.76}
\langle n|\rho(t)|m\rangle\equiv\sum_\alpha p_\alpha\langle n|\Psi_\alpha(t)\rangle\langle\Psi_\alpha(t)|m\rangle.
\end{equation}

经典地，概率（密度）$\rho(t)$ 是一个定义在相空间上的函数。像相空间中的所有算符一样，它在量子力学中被替换为一个矩阵。抛开基的选择，密度矩阵为
\begin{equation}\label{eq:6.77}
\rho(t)=\sum_\alpha p_\alpha|\Psi_\alpha(t)\rangle\langle\Psi_\alpha(t)|.
\end{equation}

显然，$\rho$ 对应于一个\emph{纯态}（\textit{pure state}）当且仅当 $\rho^2=\rho$。

密度矩阵满足如下性质：

（i）\emph{归一化}。由于每个 $|\Psi_\alpha\rangle$ 都归一化到 1，
\begin{equation}\label{eq:6.78}
\langle1\rangle=\operatorname{tr}(\rho)=\sum_n\langle n|\rho|n\rangle=\sum_{\alpha,n}p_\alpha|\langle n|\Psi_\alpha\rangle|^2=\sum_\alpha p_\alpha=1.
\end{equation}

（ii）\emph{厄米性}。密度矩阵是\emph{厄米的}（\textit{Hermitian}），即 $\rho^\dagger=\rho$，因为
\begin{equation}\label{eq:6.79}
\langle m|\rho^\dagger|n\rangle=\langle n|\rho|m\rangle^*=\sum_\alpha p_\alpha\langle\Psi_\alpha|m\rangle\langle n|\Psi_\alpha\rangle=\langle n|\rho|m\rangle,
\end{equation}
这保证了式 \eqref{eq:6.75} 中的平均值是实数。

% p.173

（iii）\emph{正定性}。对于任意 $|\Phi\rangle$，
\begin{equation}\label{eq:6.80}
\langle\Phi|\rho|\Phi\rangle=\sum_\alpha p_\alpha\langle\Phi|\Psi_\alpha\rangle\langle\Psi_\alpha|\Phi\rangle=\sum_\alpha p_\alpha|\langle\Phi|\Psi_\alpha\rangle|^2\ge0.
\end{equation}
因此，$\rho$ 是\emph{正定的}（\textit{positive definite}），这意味着它的所有本征值都必须是正的。

\term{Liouville 定理}{Liouville's theorem}{6.liouville-theorem}给出经典密度的时间演化，即
\begin{equation}\label{eq:6.81}
\frac{\mathrm d\rho}{\mathrm dt}=\frac{\partial\rho}{\partial t}-\{\mathscr H,\rho\}=0.
\end{equation}

考察量子密度矩阵的演化，最方便的做法是在能量本征态的基中进行，此时根据式 \eqref{eq:6.70} 有
\begin{equation}\label{eq:6.82}
\begin{aligned}
\mathrm i\hbar\frac{\partial}{\partial t}\langle n|\rho(t)|m\rangle&=\mathrm i\hbar\frac{\partial}{\partial t}\sum_\alpha p_\alpha\langle n|\Psi_\alpha(t)\rangle\langle\Psi_\alpha(t)|m\rangle\\
&=\sum_\alpha p_\alpha\left[(\mathcal E_n-\mathcal E_m)\langle n|\Psi_\alpha\rangle\langle\Psi_\alpha|m\rangle\right]\\
&=\langle n|(\mathscr H\rho-\rho\mathscr H)|m\rangle.
\end{aligned}
\end{equation}
最终结果是一个张量恒等式，因此与基的选择\emph{无关}，即
\begin{equation}\label{eq:6.83}
\mathrm i\hbar\frac{\partial\rho}{\partial t}=[\mathscr H,\rho].
\end{equation}

\emph{平衡}态要求平均值\emph{不随时间变化}，并提示 $\partial\rho/\partial t=0$。在式 \eqref{eq:6.81} 与 \eqref{eq:6.83} 中，通过取 $\rho=\rho(\mathscr H)$ 即可满足这一条件。（正如第 3 章所讨论的，$\rho$ 也可以依赖于守恒量 $\{L_\alpha\}$，只要 $[\mathscr H,L_\alpha]=0$。）于是，可以类比经典统计力学构造出各种平衡量子密度矩阵。

\emph{微正则系综}：由于内能具有固定值 $E$，包含这一约束的密度矩阵为
\begin{equation}\label{eq:6.84}
\rho(E)=\frac{\delta(\mathscr H-E)}{\Omega(E)}.
\end{equation}
特别地，在能量本征态的基中，
\begin{equation}\label{eq:6.85}
\langle n|\rho|m\rangle=\sum_\alpha p_\alpha\langle n|\Psi_\alpha\rangle\langle\Psi_\alpha|m\rangle=\begin{cases}\dfrac1\Omega&\text{如果 }\mathcal E_n=E\text{，而且 }m=n,\\[8pt]0&\text{如果 }\mathcal E_n\ne E\text{，或者 }m\ne n.\end{cases}
\end{equation}
第一个条件表明，量子波函数中只能出现能量正确的本征态，并且（对于 $p_\alpha=1/\mathcal N$）这类本征态平均而言具有相同的振幅，$\overline{|\langle n|\Psi\rangle|^2}=1/\Omega$。这等价于经典的\emph{等先验平衡概率假设}。第二个（附加的）条件表明，能量为 $E$ 的 $\Omega$ 个本征态以\emph{相互独立的随机相位}组合进一个典型微观态。（注意，归一化条件 $\operatorname{tr}\rho=1$ 意味着 $\Omega(E)=\sum_n\delta(E-\mathcal E_n)$ 是 $\mathscr H$ 的能量为 $E$ 的本征态数目。）

% p.174

\emph{正则系综}：固定温度 $T=1/k_B\beta$ 可以通过让系统与一个热库接触来实现。考虑复合系统，得到正则密度矩阵为
\begin{equation}\label{eq:6.86}
\rho(\beta)=\frac{\exp(-\beta\mathscr H)}{Z(\beta)}.
\end{equation}
归一化条件 $\operatorname{tr}(\rho)=1$ 给出量子配分函数
\begin{equation}\label{eq:6.87}
Z=\operatorname{tr}\left(\mathrm e^{-\beta\mathscr H}\right)=\sum_n\mathrm e^{-\beta\mathcal E_n}.
\end{equation}
最后的求和是对 $\mathscr H$ 的（分立）能级进行的，这为前几节中所做的计算提供了依据。

\emph{巨正则系综}：粒子数 $N$ 不再固定。粒子数不确定的量子微观态张成一个所谓的\term{Fock 空间}{Fock space}{6.fock-space}。相应的密度矩阵为
\begin{equation}\label{eq:6.88}
\rho(\beta,\mu)=\frac{\mathrm e^{-\beta\mathscr H+\beta\mu N}}{\mathcal Q},\qquad\text{其中}\qquad\mathcal Q(\beta,\mu)=\operatorname{tr}\left(\mathrm e^{-\beta\mathscr H+\beta\mu N}\right)=\sum_{N=0}^\infty\mathrm e^{\beta\mu N}Z_N(\beta).
\end{equation}

\textbf{例子。}考虑体积为 $V$ 的盒子中处于量子正则系综里的一个单粒子。Hamilton 量
\begin{equation}\label{eq:6.89}
\mathscr H_1=\frac{\hat{\vec p}^2}{2m}=-\frac{\hbar^2}{2m}\nabla^2\qquad\text{（在坐标基中）},
\end{equation}
由 $\mathscr H_1|\vec k\rangle=\mathcal E(\vec k)|\vec k\rangle$ 得到的能量本征态为
\begin{equation}\label{eq:6.90}
\langle\vec x|\vec k\rangle=\frac{\mathrm e^{\mathrm i\vec k\cdot\vec x}}{\sqrt V},\qquad\text{其中}\qquad\mathcal E(\vec k)=\frac{\hbar^2k^2}{2m}.
\end{equation}
在尺寸为 $L$ 的盒子中取周期性边界条件，$\vec k$ 的允许值为 $(2\pi/L)(\ell_x,\ell_y,\ell_z)$，其中 $(\ell_x,\ell_y,\ell_z)$ 是整数。这个单粒子 Hilbert 空间中的一个特定向量由它沿这些基态中每一个的分量（有无穷多个）指定。因此，量子微观态的空间比相应的 6 维经典相空间大得多。当 $L\to\infty$ 时的配分函数
\begin{equation}\label{eq:6.91}
\begin{aligned}
Z_1&=\operatorname{tr}(\rho)=\sum_{\vec k}\exp\left(-\frac{\beta\hbar^2k^2}{2m}\right)=V\int\frac{\mathrm d^3\vec k}{(2\pi)^3}\exp\left(-\frac{\beta\hbar^2k^2}{2m}\right)\\
&=\frac V{(2\pi)^3}\left(\frac{2\pi mk_BT}{\hbar^2}\right)^{3/2}=\frac V{\lambda^3},
\end{aligned}
\end{equation}

% p.175

与经典结果一致，而 $\lambda=h/\sqrt{2\pi mk_BT}$（这为把 $\mathrm d^3\vec p\,\mathrm d^3\vec q/h^3$ 用作相空间的正确无量纲测度提供了依据）。密度矩阵在坐标表象中的矩阵元为
\begin{equation}\label{eq:6.92}
\begin{aligned}
\langle\vec x'|\rho|\vec x\rangle&=\sum_{\vec k}\langle\vec x'|\vec k\rangle\frac{\mathrm e^{-\beta\mathcal E(\vec k)}}{Z_1}\langle\vec k|\vec x\rangle=\frac{\lambda^3}{V}\int\frac{V\mathrm d^3\vec k}{(2\pi)^3}\frac{\mathrm e^{-\mathrm i\vec k\cdot(\vec x-\vec x')}}{V}\exp\left(-\frac{\beta\hbar^2k^2}{2m}\right)\\
&=\frac1V\exp\left[-\frac{m(\vec x-\vec x')^2}{2\beta\hbar^2}\right]=\frac1V\exp\left[-\frac{\pi(\vec x-\vec x')^2}{\lambda^2}\right].
\end{aligned}
\end{equation}
对角元 $\langle\vec x|\rho|\vec x\rangle=1/V$ 正是在 $\vec x$ 处找到一个粒子的概率。非对角元没有经典对应。它们提示我们把粒子理解为一个尺寸为 $\lambda=h/\sqrt{2\pi mk_BT}$ 的波包，即\term{热波长}{thermal wavelength}{6.thermal-wavelength}。当 $T\to\infty$ 时，$\lambda$ 趋于零，经典分析是有效的。当 $T\to0$ 时，$\lambda$ 发散，而当 $\lambda$ 与盒子的尺寸相当时，量子力学效应开始占主导。

另外，我们也可以通过注意到式 \eqref{eq:6.86} 蕴含下式而得到式 \eqref{eq:6.92}
\begin{equation}\label{eq:6.93}
\frac{\partial}{\partial\beta}Z\rho=-\mathscr HZ\rho=\frac{\hbar^2}{2m}\nabla^2Z\rho.
\end{equation}
这正（自由粒子的）扩散方程，可以在初始条件 $\rho(\beta=0)=1$（即 $\langle\vec x'|\rho(\beta=0)|\vec x\rangle=\delta^3(\vec x-\vec x')/V$）下求解，以给出式 \eqref{eq:6.92} 的结果。

% =============================================================
%  本章斜体术语清单（供用户圈定附录 B）
% =============================================================
% —— 附录 B 候选术语（原文中的斜体术语，本章内首次出现处已埋 \term 锚点）：
%    p.157 劲度（stiffness）—— 6.stiffness
%    p.157 幺正的（unitary）—— 6.unitary
%    p.157 旋转对称性（Rotation symmetry）—— 6.rotation-symmetry
%    p.157 振动的（vibrational）—— 6.vibrational
%    p.159 量子涨落（quantum fluctuations）—— 6.quantum-fluctuations
%    p.161 晶格（lattice）—— 6.lattice
%    p.161 基矢（basis vectors）—— 6.basis-vectors
%    p.162 Fourier 基（Fourier basis）—— 6.fourier-basis
%    p.162 布里渊区（Brillouin zone）—— 6.brillouin-zone
%    p.162 波矢（wavevectors）—— 6.wavevector
%    p.164 爱因斯坦模型（Einstein model）—— 6.einstein-model
%    p.164 德拜模型（Debye model）—— 6.debye-model
%    p.165 声速（speed of sound）—— 6.speed-of-sound
%    p.166 声子（phonons）—— 6.phonons
%    p.166 色散关系（dispersion relation）—— 6.dispersion-relation
%    p.166 周期性边界条件（periodic boundary conditions）—— 6.periodic-bc
%    p.166 态密度（density of states）—— 6.density-of-states
%    p.166 德拜温度（Debye temperature）—— 6.debye-temperature
%    p.167 光子（photons）—— 6.photons
%    p.167 横模（transverse）—— 6.transverse
%    p.167 紫外灾变（ultraviolet catastrophe）—— 6.ultraviolet-catastrophe
%    p.168 Casimir 力（Casimir force）—— 6.casimir-force
%    p.169 Stefan--Boltzmann 定律（Stefan--Boltzmann law）—— 6.stefan-boltzmann-law
%    p.169 Stefan 常量（Stefan's constant）—— 6.stefan-constant
%    p.170 单位向量（unit vector）—— 6.unit-vector
%    p.170 Hilbert 空间（Hilbert space）—— 6.hilbert-space
%    p.170 波函数（wave function）—— 6.wave-function
%    p.170 厄米矩阵（Hermitian matrices）—— 6.hermitian-matrix
%    p.170 对易关系（commutation relations）—— 6.commutation-relations
%    p.171 本征态（eigenstates）—— 6.eigenstates
%    p.171 本征能量（eigenenergies）—— 6.eigenenergies
%    p.171 时间演化算符（time evolution operator）—— 6.time-evolution-operator
%    p.172 相空间密度（phase space density）—— 6.phase-space-density
%    p.172 密度矩阵（density matrix）—— 6.density-matrix
%    p.173 Liouville 定理（Liouville's theorem）—— 6.liouville-theorem
%    p.174 Fock 空间（Fock space）—— 6.fock-space
%    p.175 热波长（thermal wavelength）—— 6.thermal-wavelength
% —— 以下术语已在第 1--5 章（或本章前面）首次标记，本章此处不再重复埋锚点；
%    正文仍照原文保留其意大利斜体英文（附录 B 的页码只取全书首次出现处）：
%    p.157 正定的（positive definite，第 1 章 p.25）
%    p.157 简正模（normal modes，第 3 章 p.83）
%    p.157 正则的（canonical，第 4 章 p.99，作 canonical change of variables）
%    p.157 平移对称性（Translation symmetry，第 5 章 p.128，作 translational symmetry）
%    p.170 微观态（Microstates，第 3 章 p.57）
%    p.170 基矢（basis vectors，本章 p.161 已埋锚点，此处为同章重复出现）
%    p.172 宏观态（Macrostates，第 3 章 p.58）
%    p.172 系综（ensemble，第 3 章 p.58，作 statistical ensemble）
%    p.172 无偏的（unbiased，第 2 章 p.51，作 unbiased estimate）
%    p.172 混合态（mixed state，第 3 章 p.58）、纯态（pure state，第 3 章 p.58）
%    p.173 平衡态（Equilibrium，第 1 章 p.1）
%    p.173 微正则系综（microcanonical ensemble，第 4 章 p.98）
%    p.174 正则系综（canonical ensemble，第 4 章 p.110）、巨正则系综（grand canonical ensemble，第 4 章 p.118）
% —— 原文中属于「强调」而非术语的斜体（讲义中以 \emph{} 呈现，不收入附录 B）：
%    p.156 每个（Each）、之间（between）；p.158 量子化（quantized）、振动模式（Vibrational modes）；
%    p.160 转动模式（Rotational modes）；p.162 相对的（relative）；
%    p.164 占有数（occupation numbers，行内说明）；p.165、p.167 偏振（polarizations）；
%    p.170--171 厄米的（Hermitian）、时间演化（Time evolution）、对角化（diagonalize）、
%          无关（independent）、平均（average）；
%    p.172 归一化（Normalization）、厄米性（Hermiticity）、厄米的（Hermitian）；
%    p.173 正定性（Positivity）、无关（independent）、不随时间变化（time-independent）、
%          相互独立的随机相位（independent random phases）；
%    p.174 例子（Example，行内小标题，原文为粗体）；
%    p.175 对角元（diagonal elements）、非对角元（off-diagonal elements）。
% —— 用户拍板（第 5 章交付后）：原书小节标题与行内小标题一律为**粗体**、非斜体，
%    标题里的词不属于「原文中的斜体术语」，不收录、不埋锚点。本章小节标题
%    （6.1 Dilute polyatomic gases … 6.5 Quantum macrostates）与行内小标题
%    （如 p.174 的 \textbf{例子。}、p.158 的「振动模式：」「转动模式：」）同此处理。
% —— 用户拍板：等先验平衡概率假设（the assumption of equal a priori equilibrium
%    probabilities）**不收入附录 B**，原文斜体在讲义中以 \emph{} 呈现。
% —— 习题部分（p.175 下半页起至 p.180）不译。
